\chapter{Lumière et histogramme}
\section{Lire l’histogramme fixe}
L’histogramme reste en haut du panneau Réglages pendant le défilement.
Il décrit l’aperçu sRGB courant après traitement, pas les valeurs brutes du
capteur. Les indicateurs d’extrémité signalent des valeurs de l’aperçu à ses
bornes ; ils ne mesurent pas directement la saturation du capteur.

Les plages de lecture sont Noirs (0--10\,\%), Ombres (10--33\,\%),
Exposition (33--66\,\%), Hautes lumières (66--90\,\%) et Blancs (90--100\,\%).
Elles servent de repères visuels ; elles ne définissent pas des frontières
rigides d’action des curseurs.

\section{Les réglages}
\begin{longtable}{@{}L{3cm}L{2.8cm}L{7.8cm}@{}}
\toprule Réglage & Domaine & Utilisation\\\midrule\endhead
Exposition & $-5$ à $+5$ EV & Corriger globalement le niveau lumineux ; commencer par de petites variations.\\
Contraste & $-100$ à $+100$ & Renforcer ou adoucir la séparation des tons.\\
Ombres & $-100$ à $+100$ & Éclaircir ou assombrir les zones sombres sans cibler seulement l’extrémité noire.\\
Noirs & $-100$ à $+100$ & Ajuster les tons les plus sombres et l’assise de l’image.\\
Hautes lumières & $-100$ à $+100$ & Modérer ou renforcer les tons clairs.\\
Blancs & $-100$ à $+100$ & Ajuster les tons les plus proches du blanc.\\
\bottomrule
\end{longtable}
La valeur zéro est la référence de ces six réglages. La case \emph{Niveaux
de gris}, placée dans ce groupe, active ou désactive le rendu monochrome.

\section{Procédure conseillée}
Régler d’abord l’exposition, puis les ombres et hautes lumières, enfin les
noirs, blancs et le contraste. Revenir sur les étapes précédentes si nécessaire :
ces réglages interagissent visuellement. Une hausse des ombres rend aussi le
bruit présent plus visible. Un réglage de hautes lumières ne peut pas recréer
une information qui n’existe plus dans la source.

\paragraph{Exemple.}
Sur un contre-jour, essayer une légère baisse des hautes lumières et une
hausse modérée des ombres, sans déplacer immédiatement les deux curseurs
aux extrêmes. Maintenir \touche{Espace} dans la vue image pour contrôler
l’effet, puis vérifier les transitions au zoom avec \touche{Ctrl + 1}.
\gestesreglages
